% ========================================================
% BÀI TẬP TOÁN 5 - KẾT NỐI TRI THỨC VỚI CUỘC SỐNG
% Bài 44: Luyện tập chung
% Định dạng: BTTeX Part III - Trắc nghiệm điền số ngắn
% ========================================================

% Câu 1
\begin{ex}
  Một cửa hàng có 120 chiếc xe đạp, đã bán được 20\% số xe. Cửa hàng đã bán được bao nhiêu chiếc xe đạp?
  \shortans{24}
  \loigiai{
    Số xe đã bán là: 120 x 20% = 24 (chiếc).
  }
\end{ex}

% Câu 2
\begin{ex}
  Một cửa hàng có 140 chiếc xe đạp, đã bán được 20\% số xe. Cửa hàng đã bán được bao nhiêu chiếc xe đạp?
  \shortans{28}
  \loigiai{
    Số xe đã bán là: 140 x 20% = 28 (chiếc).
  }
\end{ex}

% Câu 3
\begin{ex}
  Một cửa hàng có 160 chiếc xe đạp, đã bán được 20\% số xe. Cửa hàng đã bán được bao nhiêu chiếc xe đạp?
  \shortans{32}
  \loigiai{
    Số xe đã bán là: 160 x 20% = 32 (chiếc).
  }
\end{ex}

% Câu 4
\begin{ex}
  Một cửa hàng có 180 chiếc xe đạp, đã bán được 20\% số xe. Cửa hàng đã bán được bao nhiêu chiếc xe đạp?
  \shortans{36}
  \loigiai{
    Số xe đã bán là: 180 x 20% = 36 (chiếc).
  }
\end{ex}

% Câu 5
\begin{ex}
  Một cửa hàng có 200 chiếc xe đạp, đã bán được 20\% số xe. Cửa hàng đã bán được bao nhiêu chiếc xe đạp?
  \shortans{40}
  \loigiai{
    Số xe đã bán là: 200 x 20% = 40 (chiếc).
  }
\end{ex}

% Câu 6
\begin{ex}
  Một cửa hàng có 220 chiếc xe đạp, đã bán được 20\% số xe. Cửa hàng đã bán được bao nhiêu chiếc xe đạp?
  \shortans{44}
  \loigiai{
    Số xe đã bán là: 220 x 20% = 44 (chiếc).
  }
\end{ex}

% Câu 7
\begin{ex}
  Một cửa hàng có 240 chiếc xe đạp, đã bán được 20\% số xe. Cửa hàng đã bán được bao nhiêu chiếc xe đạp?
  \shortans{48}
  \loigiai{
    Số xe đã bán là: 240 x 20% = 48 (chiếc).
  }
\end{ex}

% Câu 8
\begin{ex}
  Một cửa hàng có 260 chiếc xe đạp, đã bán được 20\% số xe. Cửa hàng đã bán được bao nhiêu chiếc xe đạp?
  \shortans{52}
  \loigiai{
    Số xe đã bán là: 260 x 20% = 52 (chiếc).
  }
\end{ex}

% Câu 9
\begin{ex}
  Một cửa hàng có 280 chiếc xe đạp, đã bán được 20\% số xe. Cửa hàng đã bán được bao nhiêu chiếc xe đạp?
  \shortans{56}
  \loigiai{
    Số xe đã bán là: 280 x 20% = 56 (chiếc).
  }
\end{ex}

% Câu 10
\begin{ex}
  Một cửa hàng có 300 chiếc xe đạp, đã bán được 20\% số xe. Cửa hàng đã bán được bao nhiêu chiếc xe đạp?
  \shortans{60}
  \loigiai{
    Số xe đã bán là: 300 x 20% = 60 (chiếc).
  }
\end{ex}

% Câu 11
\begin{ex}
  Một cửa hàng có 320 chiếc xe đạp, đã bán được 20\% số xe. Cửa hàng đã bán được bao nhiêu chiếc xe đạp?
  \shortans{64}
  \loigiai{
    Số xe đã bán là: 320 x 20% = 64 (chiếc).
  }
\end{ex}

% Câu 12
\begin{ex}
  Một cửa hàng có 340 chiếc xe đạp, đã bán được 20\% số xe. Cửa hàng đã bán được bao nhiêu chiếc xe đạp?
  \shortans{68}
  \loigiai{
    Số xe đã bán là: 340 x 20% = 68 (chiếc).
  }
\end{ex}

% Câu 13
\begin{ex}
  Một cửa hàng có 360 chiếc xe đạp, đã bán được 20\% số xe. Cửa hàng đã bán được bao nhiêu chiếc xe đạp?
  \shortans{72}
  \loigiai{
    Số xe đã bán là: 360 x 20% = 72 (chiếc).
  }
\end{ex}

% Câu 14
\begin{ex}
  Một cửa hàng có 380 chiếc xe đạp, đã bán được 20\% số xe. Cửa hàng đã bán được bao nhiêu chiếc xe đạp?
  \shortans{76}
  \loigiai{
    Số xe đã bán là: 380 x 20% = 76 (chiếc).
  }
\end{ex}

% Câu 15
\begin{ex}
  Một cửa hàng có 400 chiếc xe đạp, đã bán được 20\% số xe. Cửa hàng đã bán được bao nhiêu chiếc xe đạp?
  \shortans{80}
  \loigiai{
    Số xe đã bán là: 400 x 20% = 80 (chiếc).
  }
\end{ex}

% Câu 16
\begin{ex}
  Một cửa hàng có 420 chiếc xe đạp, đã bán được 20\% số xe. Cửa hàng đã bán được bao nhiêu chiếc xe đạp?
  \shortans{84}
  \loigiai{
    Số xe đã bán là: 420 x 20% = 84 (chiếc).
  }
\end{ex}

% Câu 17
\begin{ex}
  Một cửa hàng có 440 chiếc xe đạp, đã bán được 20\% số xe. Cửa hàng đã bán được bao nhiêu chiếc xe đạp?
  \shortans{88}
  \loigiai{
    Số xe đã bán là: 440 x 20% = 88 (chiếc).
  }
\end{ex}

% Câu 18
\begin{ex}
  Một cửa hàng có 460 chiếc xe đạp, đã bán được 20\% số xe. Cửa hàng đã bán được bao nhiêu chiếc xe đạp?
  \shortans{92}
  \loigiai{
    Số xe đã bán là: 460 x 20% = 92 (chiếc).
  }
\end{ex}

% Câu 19
\begin{ex}
  Một cửa hàng có 480 chiếc xe đạp, đã bán được 20\% số xe. Cửa hàng đã bán được bao nhiêu chiếc xe đạp?
  \shortans{96}
  \loigiai{
    Số xe đã bán là: 480 x 20% = 96 (chiếc).
  }
\end{ex}

% Câu 20
\begin{ex}
  Một cửa hàng có 500 chiếc xe đạp, đã bán được 20\% số xe. Cửa hàng đã bán được bao nhiêu chiếc xe đạp?
  \shortans{100}
  \loigiai{
    Số xe đã bán là: 500 x 20% = 100 (chiếc).
  }
\end{ex}

% Câu 21
\begin{ex}
  Một cửa hàng có 520 chiếc xe đạp, đã bán được 20\% số xe. Cửa hàng đã bán được bao nhiêu chiếc xe đạp?
  \shortans{104}
  \loigiai{
    Số xe đã bán là: 520 x 20% = 104 (chiếc).
  }
\end{ex}

% Câu 22
\begin{ex}
  Một cửa hàng có 540 chiếc xe đạp, đã bán được 20\% số xe. Cửa hàng đã bán được bao nhiêu chiếc xe đạp?
  \shortans{108}
  \loigiai{
    Số xe đã bán là: 540 x 20% = 108 (chiếc).
  }
\end{ex}

% Câu 23
\begin{ex}
  Một cửa hàng có 560 chiếc xe đạp, đã bán được 20\% số xe. Cửa hàng đã bán được bao nhiêu chiếc xe đạp?
  \shortans{112}
  \loigiai{
    Số xe đã bán là: 560 x 20% = 112 (chiếc).
  }
\end{ex}

% Câu 24
\begin{ex}
  Một cửa hàng có 580 chiếc xe đạp, đã bán được 20\% số xe. Cửa hàng đã bán được bao nhiêu chiếc xe đạp?
  \shortans{116}
  \loigiai{
    Số xe đã bán là: 580 x 20% = 116 (chiếc).
  }
\end{ex}

% Câu 25
\begin{ex}
  Một cửa hàng có 600 chiếc xe đạp, đã bán được 20\% số xe. Cửa hàng đã bán được bao nhiêu chiếc xe đạp?
  \shortans{120}
  \loigiai{
    Số xe đã bán là: 600 x 20% = 120 (chiếc).
  }
\end{ex}
